Dans cette partie, nous évoquerons les principaux besoins que nous avons relevé en discutant avec les chercheurs. 

\subsection{La présentation du corpus et la mise en contexte des \textit{regions}}


L'application AIKON permet de constituer un corpus avec différents \textit{witnesses} dans le but de leur faire subir des traitements automatiques dont l'extraction de \textit{regions}. A cela s'ajoute diverses métadonnées pour les décrire comme le lieu de conservation, la langue ou le type de document. Il est possible de faire des recherches par le biais d'un formulaire. Néanmoins, avec cette méthode nous nous heurtons aux problèmes liés à l'utilisation d'une barre de recherche, c'est-à-dire la nécessité d'avoir déjà au préalable une connaissance du corpus. Cela est aussi un frein à la découverte par sérendipité. Il serait donc appréciable de créer une visualisation permettant de découvrir ce corpus par différentes entrées pour arriver aux \textit{regions} des \textit{witnesses}. Il faut aussi prendre en compte l'objectif principal d'AIKON qui est l'étude de la diffusion de savoirs dans le temps et l'espace. 


\subsection{La visualisation des similarités}

L'autre principale fonctionnalité d'AIKON citée précédemment est la reconnaissance automatique de similarités. 

Les chercheurs sont amenés à annoter les résultats de la recherche automatique de similarités pour évaluer les résultats et améliorer les modèles de \textit{computer vision}. Pour l'instant, une interface permet de réaliser cette tâche mais elle se limite à une section dans chaque \textit{witness} affichant la \textit{region} ainsi que ses différentes similarités avec la possibilité de filtrer les résultats par \textit{witness} et par score de similarité. Nous n'avons donc pas de vue d'ensemble des \textit{regions} similaires au sein de deux ou plusieurs \textit{witnesses}. Créer une visualisation permettant de voir tous les résultats serait utile pour faciliter l'annotation et la comparaison.

Au delà de l'annotation, les chercheurs auraient aussi besoin de visualisations pour mener leurs recherches. Stefan Zieme est un chercheur travaillant sur les différentes traditions de l'\almageste. Lors de la conférence internationale \gls{eida} 2025\footcite{Conference2025Long2025}, il a choisi de parler de l'utilisation d'AIKON dans ses recherches sur la diffusion de la traduction de Georges de Trébizonde. Il traite ses sources avec l'extraction automatique de \textit{regions} avant de classer les différents diagrammes à la main en fonction de la période et de l'endroit dans un tableau. Pour réaliser cette tâche, il utilise la plateforme \textit{Miro} qui sert à l'origine au management de projet. Cette manière de travailler est chronophage dans la mesure où il doit placer chaque diagramme à la main. Pour l'aider dans son travail, il serait judicieux de créer une visualisation permettant d'intégrer automatiquement les différentes \textit{regions} en fonction de leur similarité avec d'autres. Cela permettrait également la mise à jour automatique de la visualisation lorsqu'il ajoute de nouvelles \textit{regions}.
Stefan Zieme aurait également besoin de cette visualisation pour présenter ses recherches dans des conférences comme celle d'\gls{eida}. Nous en venons donc à un autre point important à prendre en compte : la médiation scientifique des résultats. 

\subsection{La médiation scientifique}

Les chercheurs présentent les résultats de leurs recherches lors de cours ou de conférences. Comme nous venons de le voir avec Stefan Zieme, ils essayent déjà de créer leurs propres visualisations à l'aide de différents outils. 

Néanmoins, ce ne sont pas les seuls qui seront amenés à parcourir les résultats du projet lorsque ce dernier sera terminé et qu'une interface publique existera. Un public moins spécialiste pourra également parcourir le corpus comme c'est le cas avec \gls{dishas}.

C'est pour cette raison que des chercheurs comme Matthieu Husson, principal investigateur du projet, demande la création de visualisations ludiques et attrayantes visuellement. Ce dernier aimerait l'élaboration d'une interface similaire à celle de \textit{Coins}\footcite{COINSJourneyRich}. Il s'agit d'une \textit{generous interface} créée par l'\gls{uclab} ayant pour but l'exploration des 30000 artefacts numérisés de la collection numismatique du Munzkabinett de Berlin. Elle incite donc à se comporter comme le \og \textit{flâneur d'informations} \fg de Marian Dörk en proposant aux visiteurs de manipuler les objets de la collection, ce qui n'est pas possible dans un musée classique. Elle est caractérisée par une interactivité minimaliste. Ce que Matthieu Husson recherche lorsqu'il propose de réaliser une visualisation similaire est la possibilité d'explorer les collections sous plusieurs plans et de différentes manières. En effet, les artefacts sont disposés sous forme de tas ou de mosaïques en fonction de filtres comme le pays, la période, le poids ou le diamètre. \\

Nous allons donc établir un \textit{benchmark} des solutions envisageables afin de répondre aux besoins des chercheurs.
